\documentclass[12pt, a4paper]{article}
\usepackage[UTF8]{ctex}

\usepackage{amsmath, amssymb, amsthm}
\usepackage{geometry}
\usepackage{fancyhdr}
\usepackage{enumitem}
\usepackage{xcolor}
\usepackage{hyperref}

\geometry{margin=1in}
\pagestyle{fancy}
\fancyhf{}
\rhead{线性代数（A）}
\lhead{李生丹}
\rfoot{第 \thepage 页}

\newcommand{\problem}[1]{\subsection*{题 #1}}
\newcommand{\solution}{\paragraph{解.}}

\definecolor{solutioncolor}{RGB}{0, 100, 0}

\begin{document}

\begin{center}
    {\LARGE\bfseries 线性代数 第 3 次作业 - 自动求解} \\[0.5em]
    {\large 线性代数（A）} \\[0.3em]
    李生丹 \quad | \quad 2026-10-05
\end{center}

\hrule
\vspace{1em}
\problem{1}

计算下列矩阵的行列式：
$$A = \begin{pmatrix} 2 & 1 & 3 \\ 4 & 5 & 6 \\ 7 & 8 & 9 \end{pmatrix}$$

\solution

矩阵 $A = \left[\begin{matrix}2 & 1 & 3\\4 & 5 & 6\\7 & 8 & 9\end{matrix}\right]$

按行列式定义展开（SymPy 符号计算，逐步约化）

$\det(A) = -9$

\textbf{Answer:}
\[\boxed{-9}\]

\vspace{1em}
\hrule
\vspace{1em}

\problem{2}

求下列矩阵的逆矩阵：
$$A = \begin{pmatrix} 2 & 1 \\ 5 & 3 \end{pmatrix}$$

\solution

矩阵 $A = \left[\begin{matrix}2 & 1\\5 & 3\end{matrix}\right]$

$\det(A) = 1 \neq 0$，A 可逆

伴随矩阵法 / 初等行变换法求得 $A^{-1} = \left[\begin{matrix}3 & -1\\-5 & 2\end{matrix}\right]$

\textbf{Answer:}
\[\boxed{\left[\begin{matrix}3 & -1\\-5 & 2\end{matrix}\right]}\]

\vspace{1em}
\hrule
\vspace{1em}

\problem{3}

求矩阵的特征值与特征向量：
$$A = \begin{pmatrix} 4 & 1 \\ 2 & 3 \end{pmatrix}$$

\solution

矩阵 $A = \left[\begin{matrix}4 & 1\\2 & 3\end{matrix}\right]$

特征方程 $\det(A - \lambda I) = 0$，展开求解：

$\lambda = 5$, $\lambda = 2$

对每个特征值解 $(A - \lambda I)\xi = 0$：

$\lambda = 2$: $\xi = \left[\begin{matrix}- \frac{1}{2}\\1\end{matrix}\right]$

$\lambda = 5$: $\xi = \left[\begin{matrix}1\\1\end{matrix}\right]$

\textbf{Answer:}
\[\boxed{\lambda = 2,\ \xi = \left[\begin{matrix}- \frac{1}{2}\\1\end{matrix}\right]; \quad \lambda = 5,\ \xi = \left[\begin{matrix}1\\1\end{matrix}\right]}\]

\vspace{1em}
\hrule
\vspace{1em}

\problem{4}

解线性方程组：
$$\left[\begin{array}{ccc|c} 1 & 2 & 3 & 6 \\ 2 & 4 & 6 & 12 \\ 1 & 1 & 1 & 3 \end{array}\right]$$

\solution

增广矩阵 $\overline{A} = \left[\begin{matrix}1 & 2 & 3 & 6\\2 & 4 & 6 & 12\\1 & 1 & 1 & 3\end{matrix}\right]$

初等行变换化为行最简形，判定解的情况

$r(A) = r(\overline{A}) = 2 < n = 3$，有无穷多解（自由变量 1 个）

解集: $\left\{\left( \tau_{0}, \  3 - 2 \tau_{0}, \  \tau_{0}\right)\right\}$

\textbf{Answer:}
\[\boxed{\left\{\left( \tau_{0}, \  3 - 2 \tau_{0}, \  \tau_{0}\right)\right\}}\]

\vspace{1em}
\hrule
\vspace{1em}

\problem{5}

求下列矩阵的秩：
$$A = \begin{pmatrix} 1 & 2 & 3 \\ 2 & 4 & 6 \\ 3 & 6 & 9 \end{pmatrix}$$

\solution

矩阵 $A = \left[\begin{matrix}1 & 2 & 3\\2 & 4 & 6\\3 & 6 & 9\end{matrix}\right]$

对 A 作初等行变换化为行阶梯形

非零行数 = 1，故 $r(A) = 1$

\textbf{Answer:}
\[\boxed{1}\]

\vspace{1em}
\hrule
\vspace{1em}

\problem{6}

已知矩阵
$$A = \begin{pmatrix} 1 & 2 \\ 3 & 4 \end{pmatrix} , \quad B = \begin{pmatrix} 2 & 0 \\ 1 & 3 \end{pmatrix}$$
计算矩阵乘积 $AB$。

\solution

矩阵 $A = \left[\begin{matrix}1 & 2\\3 & 4\end{matrix}\right]$

$AB = \left[\begin{matrix}1 & 2\\3 & 4\end{matrix}\right] \cdot \left[\begin{matrix}2 & 0\\1 & 3\end{matrix}\right]$

按行乘列法则：$c_{ij} = \sum_k a_{ik} b_{kj}$

$AB = \left[\begin{matrix}4 & 6\\10 & 12\end{matrix}\right]$

\textbf{Answer:}
\[\boxed{\left[\begin{matrix}4 & 6\\10 & 12\end{matrix}\right]}\]

\vspace{1em}
\hrule
\vspace{1em}

\problem{7}

用 Gram-Schmidt 方法将下列向量组正交化：
$$\alpha_1 = \begin{pmatrix} 1 \\ 1 \\ 0 \end{pmatrix} , \quad \alpha_2 = \begin{pmatrix} 1 \\ 0 \\ 1 \end{pmatrix}$$

\solution

取 $\beta_1=\alpha_1=\begin{pmatrix}1\\1\\0\end{pmatrix}$。

计算内积 $\langle\alpha_2,\beta_1\rangle=1\cdot1+0\cdot1+1\cdot0=1$，$\langle\beta_1,\beta_1\rangle=1^2+1^2+0^2=2$。

令 $\beta_2=\alpha_2-\dfrac{\langle\alpha_2,\beta_1\rangle}{\langle\beta_1,\beta_1\rangle}\beta_1=\begin{pmatrix}1\\0\\1\end{pmatrix}-\dfrac{1}{2}\begin{pmatrix}1\\1\\0\end{pmatrix}=\begin{pmatrix}1/2\\-1/2\\1\end{pmatrix}$。

验证正交性：$\langle\beta_1,\beta_2\rangle=\dfrac12-\dfrac12+0=0$，故 $\beta_1,\beta_2$ 正交。

若需单位正交向量组，则取 $e_1=\dfrac{\beta_1}{\|\beta_1\|}=\dfrac{1}{\sqrt2}\begin{pmatrix}1\\1\\0\end{pmatrix}$，$e_2=\dfrac{\beta_2}{\|\beta_2\|}=\dfrac{1}{\sqrt{3/2}}\begin{pmatrix}1/2\\-1/2\\1\end{pmatrix}=\dfrac{1}{\sqrt6}\begin{pmatrix}1\\-1\\2\end{pmatrix}$。

\textbf{Answer:}
\[\boxed{\beta_1=\begin{pmatrix}1\\1\\0\end{pmatrix},\ \beta_2=\begin{pmatrix}1/2\\-1/2\\1\end{pmatrix};\quad e_1=\frac{1}{\sqrt2}\begin{pmatrix}1\\1\\0\end{pmatrix},\ e_2=\frac{1}{\sqrt6}\begin{pmatrix}1\\-1\\2\end{pmatrix}}\]

\vspace{1em}
\hrule
\vspace{1em}

\problem{8}

判断矩阵 $A$ 是否可对角化；若可以，求出对角化所需的矩阵 $P$ 与对角矩阵 $D$：
$$A = \begin{pmatrix} 4 & 1 \\ 2 & 3 \end{pmatrix}$$

\solution

矩阵 $A = \left[\begin{matrix}4 & 1\\2 & 3\end{matrix}\right]$

求出全部特征值与特征向量

$P = \left[\begin{matrix}-1 & 1\\2 & 1\end{matrix}\right]$, $D = \left[\begin{matrix}2 & 0\\0 & 5\end{matrix}\right]$

验证 $PDP^{-1} = A$

\textbf{Answer:}
\[\boxed{\left( \left[\begin{matrix}-1 & 1\\2 & 1\end{matrix}\right], \  \left[\begin{matrix}2 & 0\\0 & 5\end{matrix}\right]\right)}\]

\vspace{1em}
\hrule
\vspace{1em}


\vspace{2em}
\noindent\textit{Auto-generated: 8/8 problems solved by SymPy.}
\end{document}
